% GENERATED FILE. Edit canonical lessons, then run npm run generate:books.
% canonical-source-hash: fnv1a64:65bf61e6185c66c3
% canonical-lessons: ZH-C319-huidian, ZH-C319-xinxin, ZH-W319-xin-letter, ZH-C319-laidian, ZH-C319-tongxin, ZH-C319-huafei

\chapter{Calls and Confidence}
\label{ch:zh-calls-and-confidence}
\hlchaptermodality{319}

\begin{chapteropening}
\textbf{By the end of this chapter:} \emph{I can say I will call back, talk about calls coming in, keep in touch by letter and say I have confidence.}

Complete the last lesson of chapter 319: I can say I will call back, talk about calls coming in, keep in touch by letter and say I have confidence.
\end{chapteropening}

\section[\texorpdfstring{\zh{回电} (huídiàn)}{回电 (huídiàn)}]{\zh{回电} (huídiàn) — to call back; to return a call}

\label{lesson:ZH-C319-huidian}



\begin{quote}
Before the new one: say the Mandarin for a reply to a letter, then the Mandarin for to pass on a message.
\end{quote}

\subsection*{You'll want to know: \zh{回电}}
\textbf{\zh{回电}} — \emph{huídiàn} — \textquotedblleft{}to call back; to return a call\textquotedblright{}.

\begin{quote}
\textbf{\zh{我在开会一会儿给你回电}} \emph{wǒ zài kāihuì yíhuìr gěi nǐ huídiàn} — I'm in a meeting; I'll call you back in a moment.
\end{quote}

\begin{cousinweb}
\textbf{\zh{回电}} (\emph{huídiàn}) is built from \zh{回} 'to turn back' and \zh{电} 'electricity', both inherited from Middle Chinese, \emph{hwoj} and \emph{denH} (Guangyun, in Baxter's transcription), which glosses the second 'lightning'.
\end{cousinweb}

\begin{grammarlens}[title={What you've built}]
One more word for telling what happened.
\end{grammarlens}

\subsection*{Your turn}
\begin{itemize}
\raggedright
  \item \emph{Say it:} \emph{huídiàn}
  \item \emph{Say it:} \emph{huídiàn}, once more
  \item \emph{Say it:} \emph{wǒ zài kāihuì yíhuìr gěi nǐ huídiàn} — I'm in a meeting; I'll call you back in a moment.
  \item \emph{Recall:} say the Mandarin for a reply to a letter, then the Mandarin for to pass on a message, then say \emph{huídiàn} again
\end{itemize}

\begin{tcolorbox}[breakable,colback=teal!4,colframe=teal!35!black,title={Before you move on}]
What does \zh{回电} mean? (\textquotedblleft{}To call back; to return a call\textquotedblright{}.) Say it once more.
\end{tcolorbox}
\section[\texorpdfstring{\zh{信心} (xìnxīn)}{信心 (xìnxīn)}]{\zh{信心} (xìnxīn) — confidence; faith}

\label{lesson:ZH-C319-xinxin}



\begin{quote}
Before the new one: say the Mandarin for to pass on a message, then the Mandarin for to return a call.
\end{quote}

\subsection*{You'll want to know: \zh{信心}}
\textbf{\zh{信心}} — \emph{xìnxīn} — \textquotedblleft{}confidence; faith\textquotedblright{}.

\begin{quote}
\textbf{\zh{我对你有信心}} \emph{wǒ duì nǐ yǒu xìnxīn} — I have faith in you.
\end{quote}

The new character, in pieces: \zh{信} — nine strokes: the person piece \zh{亻} on the left; on the right a dot, three bars and \zh{口}. The next lesson puts it in your hand.

\begin{cousinweb}
\textbf{\zh{信心}} (\emph{xìnxīn}) is built from \zh{信} 'to trust' and \zh{心} 'the heart', both inherited from Middle Chinese, \emph{sinH} and \emph{sim} (Guangyun, in Baxter's transcription), which glosses the first 'loyal and true' and 'sincere'.
\end{cousinweb}

\begin{grammarlens}[title={What you've built}]
One more word for telling what happened.
\end{grammarlens}

\subsection*{Your turn}
\begin{itemize}
\raggedright
  \item \emph{Say it:} \emph{xìnxīn}
  \item \emph{Say it:} \emph{xìnxīn}, once more
  \item \emph{Say it:} \emph{wǒ duì nǐ yǒu xìnxīn} — I have faith in you.
  \item \emph{Recall:} say the Mandarin for to pass on a message, then the Mandarin for to return a call, then say \emph{xìnxīn} again
\end{itemize}

\begin{tcolorbox}[breakable,colback=teal!4,colframe=teal!35!black,title={Before you move on}]
What does \zh{信心} mean? (\textquotedblleft{}Confidence; faith\textquotedblright{}.) Say it once more.
\end{tcolorbox}
\section[\texorpdfstring{\zh{信} (xìn)}{信 (xìn)}]{\zh{信} — the person piece, then a dot, three bars and \zh{口}}

\label{lesson:ZH-W319-xin-letter}



\begin{quote}
The word first, then the character inside it.

\begin{itemize}
\raggedright
  \item \emph{Recall:} say \emph{xìnxīn} — \textbf{R1}, one lesson back, before it decays
\end{itemize}

\begin{quote}
\zh{你} \zh{说}
\end{quote}

\emph{Read it:} each one aloud, with nothing said first — \emph{nǐ}, \emph{shuō}
\end{quote}

\subsection*{Script — the person piece, then a dot, three bars and \zh{口}}
\hlblockfigure{\detokenize{figures/ZH-W319-xin-letter-filmstrip.pdf}}{How \zh{信} is written, stroke by stroke}

\begin{quote}
\zh{信}
\end{quote}

\textbf{Two pieces, left before right.} Count the pieces, not the strokes.

\textbf{The left piece, \zh{亻}, two strokes.} A stroke falling to the left, then a vertical down to the foot, as in \zh{你}.

\textbf{The right piece, seven strokes, top to bottom.} A dot. Three bars, the top one the longest. Then \zh{口} at the foot: a short vertical, a bar that runs right and turns down, and the bar that closes it.

2 + 7 = nine.

\subsection*{Your turn}
\begin{itemize}
\raggedright
  \item \emph{Write it:} \zh{信}, all nine, with the model in front of you
  \item \emph{Cover:} write it again from memory, and count the pieces, not the strokes
\end{itemize}

\begin{tcolorbox}[breakable,colback=teal!4,colframe=teal!35!black,title={Before you move on}]
How many strokes in \zh{信}? (\textbf{Nine}.) What is the left piece? (\textbf{\zh{亻}, the person piece}.) What closes the right piece at the foot? (\textbf{\zh{口}}.)

Source: \href{https://raw.githubusercontent.com/chanind/hanzi-writer-data/68d10a4b21150cae5e1ebbd223eed289cf32d90c/data/\%E4\%BF\%A1.json}{Hanzi Writer Data \zh{信}, pinned stroke paths and medians}.
\end{tcolorbox}
\section[\texorpdfstring{\zh{来电} (láidiàn)}{来电 (láidiàn)}]{\zh{来电} (láidiàn) — incoming call}

\label{lesson:ZH-C319-laidian}



\begin{quote}
Before the new one: say the Mandarin for to return a call, then the Mandarin for confidence that it will work.
\end{quote}

\subsection*{You'll want to know: \zh{来电}}
\textbf{\zh{来电}} — \emph{láidiàn} — \textquotedblleft{}incoming call\textquotedblright{}.

\begin{quote}
\textbf{\zh{开会的时候我有三个来电}} \emph{kāihuì de shíhou wǒ yǒu sān ge láidiàn} — During the meeting three calls came in for me.
\end{quote}

\begin{cousinweb}
\textbf{\zh{来电}} (\emph{láidiàn}) is built from \zh{来} 'to come' and \zh{电} 'electricity', both inherited from Middle Chinese, \emph{loj} and \emph{denH} (Guangyun, in Baxter's transcription), which glosses the first 'to arrive' and the second 'lightning'.
\end{cousinweb}

\begin{grammarlens}[title={What you've built}]
One more word for telling what happened.
\end{grammarlens}

\subsection*{Your turn}
\begin{itemize}
\raggedright
  \item \emph{Say it:} \emph{láidiàn}
  \item \emph{Say it:} \emph{láidiàn}, once more
  \item \emph{Say it:} \emph{kāihuì de shíhou wǒ yǒu sān ge láidiàn} — During the meeting three calls came in for me.
  \item \emph{Recall:} say the Mandarin for to return a call, then the Mandarin for confidence that it will work, then say \emph{láidiàn} again
\end{itemize}

\begin{tcolorbox}[breakable,colback=teal!4,colframe=teal!35!black,title={Before you move on}]
What does \zh{来电} mean? (\textquotedblleft{}Incoming call\textquotedblright{}.) Say it once more.
\end{tcolorbox}
\section[tōngxìn]{tōngxìn — to write to each other; to correspond}

\label{lesson:ZH-C319-tongxin}



\begin{quote}
Before the new one: say the Mandarin for confidence that it will work, then the Mandarin for an incoming call.
\end{quote}

\subsection*{You'll want to know: tōngxìn}
\emph{tōngxìn} — \textquotedblleft{}to write to each other; to correspond\textquotedblright{}.

\begin{quote}
\emph{wǒ hé wǒ de bǐyǒu tōngxìn le wǔ nián} — My pen pal and I wrote to each other for five years.
\end{quote}

Its characters come later; for now it is a word for the ear.

\begin{cousinweb}
\emph{tōngxìn} is built from \emph{tōng} 'to get through' and \zh{信} 'a letter', both inherited from Middle Chinese, \emph{thuwng} and \emph{sinH} (Guangyun, in Baxter's transcription), which glosses the first 'to reach' and the second 'true'.
\end{cousinweb}

\begin{grammarlens}[title={What you've built}]
One more word for telling what happened.
\end{grammarlens}

\subsection*{Your turn}
\begin{itemize}
\raggedright
  \item \emph{Say it:} \emph{tōngxìn}
  \item \emph{Say it:} \emph{tōngxìn}, once more
  \item \emph{Say it:} \emph{wǒ hé wǒ de bǐyǒu tōngxìn le wǔ nián} — My pen pal and I wrote to each other for five years.
  \item \emph{Recall:} say the Mandarin for confidence that it will work, then the Mandarin for an incoming call, then say \emph{tōngxìn} again
\end{itemize}

\begin{tcolorbox}[breakable,colback=teal!4,colframe=teal!35!black,title={Before you move on}]
What does tōngxìn mean? (\textquotedblleft{}To write to each other; to correspond\textquotedblright{}.) Say it once more.
\end{tcolorbox}
\section[huàfèi]{huàfèi — phone charges}

\label{lesson:ZH-C319-huafei}



\begin{quote}
Before the new one: say the Mandarin for an incoming call, then the Mandarin for to write to each other.
\end{quote}

\subsection*{You'll want to know: huàfèi}
\emph{huàfèi} — \textquotedblleft{}phone charges\textquotedblright{}.

\begin{quote}
\emph{zhège yuè huàfèi yìbǎi kuài} — This month's phone bill is a hundred yuan.
\end{quote}

Its characters come later; for now it is a word for the ear.

\begin{cousinweb}
\emph{huàfèi} is built from \zh{话} 'talk' and \emph{fèi} 'a cost', both inherited from Middle Chinese, \emph{hwaejH} and \emph{phj+jH} (Guangyun, in Baxter's transcription), which glosses the second 'to use up'.
\end{cousinweb}

\begin{grammarlens}[title={What you've built}]
One more word for telling what happened.
\end{grammarlens}

\subsection*{Your turn}
\begin{itemize}
\raggedright
  \item \emph{Say it:} \emph{huàfèi}
  \item \emph{Say it:} \emph{huàfèi}, once more
  \item \emph{Say it:} \emph{zhège yuè huàfèi yìbǎi kuài} — This month's phone bill is a hundred yuan.
  \item \emph{Recall:} say the Mandarin for an incoming call, then the Mandarin for to write to each other, then say \emph{huàfèi} again
\end{itemize}

\begin{tcolorbox}[breakable,colback=teal!4,colframe=teal!35!black,title={Before you move on}]
What does huàfèi mean? (\textquotedblleft{}Phone charges\textquotedblright{}.) Say it once more.
\end{tcolorbox}
